\section{Theoretical Background and Clinical Concepts}

\subsection{Clinical Deployment Requirements}
In epilepsy monitoring, sensitivity (detecting all
seizures to ensure safety) must be weighed against precision
(minimizing false alarms). In a clinical ward, high sensitivity is
prioritized as medical staff can verify alarms. In ambulatory (home)
settings, however, a high FAR/h causes "alarm
fatigue," rendering the device useless for the patient.

Finally, the "Black Box" nature of deep learning poses safety
challenges. Unlike decision trees or logistic regression, deep neural networks
often lack transparency. For clinical acceptance, systems generally
require explainability to support physician judgment and meet regulatory requirements.

\subsection{Seizure Classification and Clinical Implications}

Epileptic seizures are classified by their onset as focal (starting in a specific area of the brain)
or generalized (affecting both hemispheres from the onset).

Focal seizures are further subdivided into focal aware and focal impaired awareness seizures,
depending on whether the patient is aware during the seizure or not, and by onset features (motor vs non‑motor).
Focal seizures may evolve to bilateral tonic–clonic seizures.

Generalized seizures begin in both cerebral hemispheres and include motor
(e.g. tonic–clonic, myoclonic, atonic) and non‑motor (absence) types
\parencite{fisherInstructionManualILAE2017}. 
Tonic–clonic seizures commonly cause loss of consciousness and carry increased
injury and SUDEP (Sudden Unexpected Death in Epilepsy) risk \parencite{hesdorfferCombinedAnalysisRisk2011}.

\subsection{Autonomic Markers and Wearable Sensing}
Sympathetic activation around seizures manifests as ictal tachycardia,
preictal HR rise with reduced HRV (lower RMSSD/HF, higher LF/HF),
increased EDA (tonic SCL, phasic SCRs), and limited ACC changes
useful for artifacts/context. 
PPG yields pulse‑rate variability analogous to HRV from ECG
\parencite{mironAutonomicBiosignalsSeizure2025,masonHeartRateVariability2024}. 
Wearable heart‑rate sensing aligns closely with chest‑strap ECG under motion (e.g., Apple Watch
Ultra/Polar H10 concordance $\sim$0.998), though dominant‑wrist
motion degrades accuracy. Chest‑strap ECG provides robust RR intervals
during exercise \parencite{wolfHearttoWearAssessingAccuracy2024a}.

\subsection{Study Phases: Retrospective vs. Pseudo-Prospective vs. Prospective}
Retrospective studies use pre-recorded data, giving the model access to the full dataset including future samples. 
Pseudo-prospective designs simulate real-time performance by processing data sequentially without future knowledge, 
providing a more realistic performance estimate. Prospective studies collect and 
process data in true real-time, representing the highest standard for clinical validity and 
the scenario in which alarm-based metrics (FAR/h) matter most.
